\section{Discussion}
\label{sec:discussion}

\paragraph{Two signals, one lever, and none of them is accommodation.}
The results separate three things that ``an internal state transition'' conflates. \emph{Sufficiency} is a state of the input that the model represents cleanly and a text model can also compute; its causal counterpart is an abstention gate. \emph{Uptake} is a state of the model: whether the same decisive paragraph will be integrated is predictable from activations and not from text, at $0.78$ AUROC across three families and on every external source. Neither readout direction is a lever; what turns an insufficient state into an answer is high-rank, orthogonal to both, and on some models concentrated on the answer's unembedding direction. The picture is a low-dimensional gate (answer or abstain) plus high-dimensional answer content, not a shared accommodation direction.

\paragraph{Why probes do not find the lever, and what to use.}
A logistic probe maximises separability along directions of low within-class variance; the mean-difference direction has a large component along a high-variance ``answering'' axis that the discriminative weights down-weight, and the projection-out and rank experiments show the causal content is exactly there. A high-AUROC, transferable probe direction can thus be causally irrelevant, which argues for validating any proposed steering direction with the projection-out, rank and matched-random controls used here. For retrieval systems the usable signal is uptake, as a detector of ignored evidence; sufficiency is tracked about as well by the model's own abstention, and the abstention gate trades abstention against compliance uniformly rather than making the model more evidence-sensitive.

\paragraph{Limitations.}
Causal experiments on the external sources use Qwen only; RAGBench correctness relies on an LLM judge; Wikipedia-trained directions transfer only partially to other domains; uptake weakens on 3--4-hop document-disjoint splits ($0.61$); and only additive and projection interventions at the anchor were tested.
